\chapter{Esercizi: Ricorsione}
\label{chap:esercizi-ricorsione}

\begin{exercise}
\label{ex:fi-esercizi-ricorsione-significato}
Che cosa si intende per ricorsione in programmazione?
\end{exercise}

La \emph{ricorsione} è una tecnica che risolve un problema mediante
sottoproblemi della stessa natura, utilizzando chiamate a funzioni
che possono invocare se stesse. Il procedimento distingue un
\emph{caso base} e un \emph{passo ricorsivo}
(\cref{sec:fi-ricorsione-terminazione}).

\begin{exercise}
\label{ex:fi-esercizi-ricorsione-funzione}
Che cosa significa dire che una funzione è ricorsiva?
\end{exercise}

Una funzione è \emph{ricorsiva} se, durante la propria esecuzione,
può invocare se stessa.
Ogni chiamata costituisce una nuova invocazione con parametri propri
(\cref{sec:fi-ricorsione-terminazione}).

\begin{exercise}
\label{ex:fi-esercizi-ricorsione-chiamata}
Quando una funzione ricorsiva chiama se stessa, che cosa succede
all'esecuzione della funzione chiamante?
\end{exercise}

L'invocazione chiamante rimane sospesa finché quella chiamata non
termina; il suo stato viene conservato in un \emph{frame} nella
pila delle chiamate. Al ritorno della chiamata interna, l'esecuzione
riprende dal punto in cui era stata sospesa
(\cref{sec:fi-pila-ricorsione}).

\begin{exercise}
\label{ex:fi-esercizi-ricorsione-fattoriale-base}
Nel calcolo del fattoriale, qual è il caso base?
\end{exercise}

Il caso base è \(n=0\), per il quale \(0!=1\): la funzione
restituisce \(1\) senza effettuare altre chiamate ricorsive
(\cref{eq:fi-fattoriale,sec:fi-fattoriale}).

\begin{exercise}
\label{ex:fi-esercizi-ricorsione-fattoriale-passo}
Nel calcolo del fattoriale, qual è il passo ricorsivo?
\end{exercise}

Per \(n>0\) si applica la relazione
\[
  n!=n(n-1)!,
\]
calcolando il fattoriale di \(n-1\) e moltiplicandone il risultato
per \(n\) (\cref{eq:fi-fattoriale}).

\begin{exercise}
\label{ex:fi-esercizi-ricorsione-fattoriale-semplice}
Perché nel fattoriale ricorsivo la chiamata
\lstinline|fattoriale(n - 1)| rappresenta un problema più semplice?
\end{exercise}

Per \(n>0\), l'argomento \(n-1\) è minore di \(n\) e resta
non negativo: il sottoproblema richiede lo stesso calcolo su un
ingresso più piccolo. Ripetendo la riduzione, si raggiunge
\(n=0\), cioè il caso base
(\cref{sec:fi-fattoriale,sec:fi-ricorsione-terminazione}).

\begin{exercise}
\label{ex:fi-esercizi-ricorsione-senza-base}
Che cosa potrebbe succedere se in una funzione ricorsiva mancasse
il caso base?
\end{exercise}

Le chiamate potrebbero continuare senza raggiungere una condizione
di arresto. In CPython, l'accumulo dei frame oltre il limite di
profondità provoca generalmente \texttt{RecursionError}; ciò non
costituisce una terminazione corretta dell'algoritmo
(\cref{sec:fi-ricorsione-terminazione,sec:fi-ricorsione-coda}).

\begin{exercise}
\label{ex:fi-esercizi-ricorsione-base-insufficiente}
Perché non basta avere un caso base: perché è importante anche che
ogni chiamata si avvicini al caso base?
\end{exercise}

Un caso base garantisce l'arresto soltanto se viene raggiunto per
ogni ingresso ammissibile. Occorre quindi che le chiamate
riducano progressivamente il problema; un criterio sufficiente è
una misura intera non negativa che diminuisca strettamente a ogni
chiamata ricorsiva (\cref{sec:fi-ricorsione-terminazione}).

\begin{exercise}
\label{ex:fi-esercizi-ricorsione-terminazione-regole}
Quali sono le due regole fondamentali che permettono di dire che
un algoritmo ricorsivo termina?
\end{exercise}

Devono essere previsti casi base che non effettuino altre chiamate
ricorsive e ogni catena di chiamate, per ogni ingresso ammissibile,
deve raggiungerne uno. Per dimostrarlo si può individuare una misura
intera non negativa strettamente decrescente, assumendo che ogni
invocazione svolga un numero finito di operazioni e generi un numero
finito di chiamate (\cref{sec:fi-ricorsione-terminazione}).

\begin{exercise}
\label{ex:fi-esercizi-ricorsione-traccia-quattro}
Calcola a mano le chiamate necessarie per eseguire
\texttt{fattoriale(4)} usando la versione ricorsiva.
\end{exercise}

Le cinque invocazioni avvengono nell'ordine
\[
\begin{aligned}
  \texttt{fattoriale(4)}&\to\texttt{fattoriale(3)}
  \to\texttt{fattoriale(2)}\\
  &\to\texttt{fattoriale(1)}\to\texttt{fattoriale(0)}.
\end{aligned}
\]
Raggiunto il caso base, i risultati vengono restituiti in ordine
inverso:
\[
\begin{aligned}
  \texttt{fattoriale(0)}&=1,\\
  \texttt{fattoriale(1)}&=1\cdot1=1,\\
  \texttt{fattoriale(2)}&=2\cdot1=2,\\
  \texttt{fattoriale(3)}&=3\cdot2=6,\\
  \texttt{fattoriale(4)}&=4\cdot6=24.
\end{aligned}
\]
La sospensione e la ripresa delle invocazioni sono illustrate
nella \cref{fig:fi-fattoriale-chiamate}.

\begin{exercise}
\label{ex:fi-esercizi-ricorsione-leggibilita}
In quali casi la ricorsione può rendere il codice più leggibile?
\end{exercise}

Quando il problema si esprime naturalmente in termini di
sottoproblemi della stessa natura, come nel fattoriale o in un
algoritmo \emph{divide et impera}, la funzione ricorsiva può
seguire direttamente la definizione del procedimento. La scelta
resta subordinata alla chiarezza dei casi base e dei passi
ricorsivi (\cref{sec:fi-fattoriale,sec:fi-ordinamento-fusione}).

\begin{exercise}
\label{ex:fi-esercizi-ricorsione-efficienza}
Perché una funzione ricorsiva può essere meno efficiente di un
ciclo equivalente?
\end{exercise}

Ogni chiamata comporta un costo di invocazione e mantiene un
frame nella pila: per il fattoriale, la versione ricorsiva richiede
spazio ausiliario \(\Theta(n)\), mentre quella iterativa richiede
\(\Theta(1)\), nel modello di costo adottato. Alcune ricorsioni
ricalcolano inoltre sottoproblemi già risolti, come nel caso di
Fibonacci (\cref{sec:fi-pila-ricorsione,sec:fi-fibonacci}).

\begin{exercise}
\label{ex:fi-esercizi-ricorsione-coda}
Che cosa si intende per ricorsione in coda?
\end{exercise}

Una chiamata ricorsiva è \emph{in coda} quando la funzione chiamante
non deve eseguire ulteriori operazioni dopo il suo ritorno. In
particolare, \lstinline|return fattoriale_coda(n - 1, acc * n)|
è in coda, mentre \lstinline|return n * fattoriale(n - 1)|
non lo è (\cref{sec:fi-ricorsione-coda}).

\begin{exercise}
\label{ex:fi-esercizi-ricorsione-coda-ciclo}
Perché la ricorsione in coda può essere trasformata abbastanza
facilmente in un ciclo?
\end{exercise}

Poiché dopo la chiamata ricorsiva non restano operazioni sospese,
lo stato necessario al passo successivo può essere aggiornato
nelle variabili di un ciclo, anziché passato a una nuova
invocazione. Questa trasformazione è mostrata per il fattoriale
nella \cref{sec:fi-ricorsione-coda}; CPython non la esegue
automaticamente.

\begin{exercise}
\label{ex:fi-esercizi-ricorsione-multipla}
Che cosa si intende per ricorsione multipla?
\end{exercise}

Si ha \emph{ricorsione multipla} quando una stessa invocazione può
eseguire più chiamate ricorsive. Per esempio, \texttt{fib(n)}
richiama \texttt{fib(n - 1)} e \texttt{fib(n - 2)} per \(n>2\)
(\cref{sec:fi-fibonacci}).

\begin{exercise}
\label{ex:fi-esercizi-ricorsione-fibonacci-inefficiente}
Perché l'algoritmo ricorsivo per Fibonacci è molto inefficiente?
\end{exercise}

La traduzione diretta della ricorrenza ricalcola più volte gli stessi
termini, producendo un numero esponenziale di invocazioni. Nel
modello a costo unitario, il tempo è \(\Theta(F_n)\), mentre la
versione iterativa richiede \(\Theta(n)\) operazioni
(\cref{sec:fi-fibonacci,sec:fi-complessita-fibonacci}).

\begin{exercise}
\label{ex:fi-esercizi-ricorsione-fibonacci-albero}
Che cosa rappresenta l'albero delle invocazioni di Fibonacci?
\end{exercise}

Ogni nodo rappresenta un'invocazione di \texttt{fib} e i suoi
figli rappresentano le chiamate ricorsive che essa effettua.
L'albero mostra tutte le invocazioni, comprese quelle ripetute,
non soltanto quelle contemporaneamente presenti nella pila.
Per \texttt{fib(5)} comprende nove invocazioni
(\cref{fig:fi-fibonacci-albero,sec:fi-pila-ricorsione}).

\begin{exercise}
\label{ex:fi-esercizi-ricorsione-somma-naturali}
Scrivi una funzione Python ricorsiva che calcoli la somma dei primi
\(n\) numeri naturali, cioè \(1+2+\cdots+n\). Indica chiaramente
il caso base e il passo ricorsivo.
\end{exercise}

\begin{unipdcode}[language=Python]
def somma_naturali(n):
    if n < 0:
        raise ValueError("n deve essere non negativo")
    if n == 0:
        return 0
    return n + somma_naturali(n - 1)
\end{unipdcode}

Per \(n=0\) la somma è nulla (\emph{caso base}); per \(n>0\)
la funzione restituisce \(n\) più la somma dei primi \(n-1\)
naturali (\emph{passo ricorsivo}). L'argomento diminuisce fino a
zero, garantendo la terminazione per gli interi non negativi
(\cref{sec:fi-esempi-ricorsione}).

\begin{exercise}
\label{ex:fi-esercizi-ricorsione-conto-rovescia}
Scrivi una funzione Python ricorsiva che stampi un conto alla
rovescia da \(n\) fino a \(0\). Spiega perché, a ogni chiamata,
il problema diventa più semplice.
\end{exercise}

\begin{unipdcode}[language=Python]
def conto_rovescia(n):
    if n < 0:
        raise ValueError("n deve essere non negativo")
    print(n)
    if n == 0:
        return
    return conto_rovescia(n - 1)
\end{unipdcode}

La funzione stampa \(n\) prima di richiamare se stessa con
\(n-1\). Per ogni intero \(n\ge0\), l'argomento diminuisce di
uno fino al caso base \(n=0\), che stampa zero senza altre chiamate.
Per \(n=3\) l'output è \(3,2,1,0\), un valore per riga
(\cref{sec:fi-esempi-ricorsione}).
